\documentclass[11pt,a4paper]{article}

\usepackage[utf8]{inputenc}
\usepackage[T1]{fontenc}
\usepackage[turkish]{babel}
\usepackage[margin=2.4cm]{geometry}
\setlength{\headheight}{14pt}
\usepackage{graphicx}
\usepackage{float}
\usepackage{booktabs}
\usepackage{tabularx}
\usepackage{amsmath,amssymb}
\usepackage{gensymb}
\usepackage[hidelinks]{hyperref}
\usepackage{fancyhdr}
\usepackage{enumitem}
\usepackage{xcolor}

\definecolor{calibblue}{rgb}{0.08,0.22,0.58}
\definecolor{datagreen}{rgb}{0.10,0.42,0.22}
\definecolor{replaypurple}{rgb}{0.38,0.18,0.62}
\definecolor{warnorange}{rgb}{0.72,0.32,0.08}

\pagestyle{fancy}
\fancyhf{}
\fancyhead[L]{\small Hafta 5 -- Kalibrasyon, Veri Dogrulama, Replay ve Analiz}
\fancyhead[R]{\small \thepage}
\renewcommand{\headrulewidth}{0.4pt}

\begin{document}

\begin{center}
    {\LARGE\bfseries Hafta 5 Geliştirme Raporu:\\[4pt]
    \textcolor{calibblue}{Zorunlu Kalibrasyon Akışı},\\[4pt]
    \textcolor{datagreen}{Veri Doğrulama ve Kayıt Aktarımı},\\[4pt]
    \textcolor{replaypurple}{Replay Sistemi} ve\\[4pt]
    \textcolor{warnorange}{Detaylı Analiz Araçları}}\\[12pt]
    {\normalsize Mayıs 2026}
\end{center}

\vspace{6pt}

\begin{abstract}
Bu rapor, proje geliştirme sürecinin beşinci haftasında tamamlanan
kalibrasyon güvenliği, veri doğrulama, kayıt saklama, replay ve analiz
araçları çalışmalarını belgelemektedir. Hafta~4'te sistem; ghost avatar
demo gösterimleri, üst ekstremite takip katmanları ve tez taslağı ile
kullanıcıya gösterilebilir ve akademik olarak raporlanabilir hale
getirilmişti. Hafta~5'te ise odak, bu yapının gerçek build ve headset
testlerine hazırlanmasına kaymıştır. Bu kapsamda kalibrasyon yapılmadan
oyunun başlatılması engellenmiş, sonuçların doğru yorumlanması için
JSON/CSV raporları ve diagnostics kayıtları incelenmiş, örnek oyun
oturumlarından elde edilen veriler bilgisayara aktarılmış, tam oturum
replay sistemi eklenmiş ve kopyalanan kayıtların ayrıntılı analizini
sağlayan Unity Editor araçları geliştirilmiştir.
\end{abstract}

\tableofcontents
\vspace{8pt}

% ================================================================
\section{Rapor Kapsamı ve Hafta 4 Üzerine Eklenenler}
\label{sec:kapsam}

Hafta~4 raporu üç ana teslim etrafında yapılandırılmıştı: görev bazlı
örnek hareketleri gösteren ghost avatar demo sistemi, daha kararlı üst
ekstremite takip katmanı ve bitirme tezi taslağı. Bu çalışmalar sistemin
kullanıcıya anlatılabilir, görsel olarak yönlendirilebilir ve akademik
biçimde raporlanabilir hale gelmesini sağlamıştı.

Hafta~5'te bu temel üzerine eklenen ana fikir şudur: sistem artık yalnızca
sahnede çalışan bir prototip değil, \textbf{kalibrasyon zorunluluğu olan},
\textbf{oturum sonunda veri saklayan}, \textbf{headset'ten bilgisayara
aktarılabilen}, \textbf{Unity içinde tekrar oynatılabilen} ve
\textbf{ayrıntılı analiz edilebilen} bir araştırma aracına dönüştürülmüştür.

\begin{table}[H]
\centering
\caption{Hafta 4 ve Hafta 5 odak farkı}
\label{tab:hafta45}
\renewcommand{\arraystretch}{1.25}
\begin{tabularx}{\textwidth}{lXX}
\toprule
\textbf{Başlık} & \textbf{Hafta 4} & \textbf{Hafta 5} \\
\midrule
Kullanıcı yönlendirme & Ghost avatar ile örnek hareket gösterimi & Kalibrasyon tamamlanmadan oturum başlatmayı engelleyen akış \\
Veri üretimi & Görev motoru ve skor hesaplama altyapısı & JSON/CSV rapor, diagnostics CSV ve replay veri paketleri \\
Test yaklaşımı & Editor ve sahne kurulum odaklı doğrulama & Build/headset örnek oturumları ve kopyalanan kayıtların incelenmesi \\
Sonradan inceleme & Tez taslağı ve raporlanabilir mimari & Replay oynatma, timeline, frame analizi ve toplu rapor export \\
\bottomrule
\end{tabularx}
\end{table}

Bu haftanın çalışmaları beş somut başlıkta toplanabilir:

\begin{enumerate}[nosep]
    \item Kalibrasyon yapılmadan oyunun başlatılamaması,
    \item Sonuç doğruluğu için rapor ve diagnostics testleri,
    \item Örnek oyun oturumlarının oynanması ve verilerin \texttt{KayitSonuclari} klasörüne aktarılması,
    \item Tam oturum replay kayıt ve oynatma sisteminin eklenmesi,
    \item Kayıtların ayrıntılı incelenmesi için Unity Editor analiz araçlarının yazılması.
\end{enumerate}

% ================================================================
\section{Kalibrasyon Yapılmadan Oyunun Başlatılamaz Hale Getirilmesi}
\label{sec:kalibrasyon}

Hafta~5'in ilk kritik düzeltmesi, oyunun kalibrasyon yapılmadan
başlayabilmesi problemine yöneliktir. Önceki geliştirme aşamalarında editor
testlerini kolaylaştırmak için simülasyon modu ve otomatik başlatma
seçenekleri bulunmaktaydı. Ancak gerçek headset/build testlerinde bu durum
araştırma verisinin güvenilirliğini zayıflatır; çünkü görev sonuçları
kalibre edilmemiş tracker-avatar eşleşmesine dayanabilir.

Bu nedenle aktif hardware sahnesinde gerçek test akışı şu şekilde
kilitlenmiştir:

\begin{center}
\textit{Tracker bekleme $\rightarrow$ kalibrasyon yönergesi $\rightarrow$
kalibrasyon tutma $\rightarrow$ kalibrasyon tamamlandı $\rightarrow$
kullanıcı başlatır $\rightarrow$ oturum ve kayıt başlar}
\end{center}

\subsection{Session Başlatma Öncesi Validasyon}

\texttt{GameFlowController.StartSession()} artık yalnızca
\texttt{FullBodyCalibrator.IsCalibrated} durumuna bakmakla kalmaz; gerçek
oturum için bir dizi hazırlık koşulunu da kontrol eder. Bu kontrol,
kalibrasyon kapısının yalnızca görsel bir ekran değil, gerçek bir çalışma
zamanı güvenlik bariyeri olmasını sağlar.

\begin{table}[H]
\centering
\caption{Oturum başlatma öncesi kontrol edilen koşullar}
\label{tab:session_gate}
\renewcommand{\arraystretch}{1.22}
\begin{tabularx}{\textwidth}{lX}
\toprule
\textbf{Koşul} & \textbf{Amaç} \\
\midrule
\texttt{FullBodyCalibrator.IsCalibrated} & Kullanıcının T-poz kalibrasyonunu tamamladığını doğrulamak \\
\texttt{FullBodyIKSolver.IsCalibrated} & Avatar IK kalibrasyonunun geçerli olduğunu doğrulamak \\
\texttt{FullBodyTrackingManager.IsAssigned} & Trackerların pelvis/ayak rollerine atanmış olduğunu doğrulamak \\
\texttt{LowerLimbBiometrics.useSimulatedInput = false} & Gerçek headset testinde simüle veri kullanılmasını engellemek \\
\texttt{SessionReportWriter} & Görev sonuçlarının oturum sonunda dosyaya yazılmasını garanti etmek \\
\texttt{ReplayRecorder} & Oyun başladığında replay kaydının alınmasını garanti etmek \\
\bottomrule
\end{tabularx}
\end{table}

Bu yapı sayesinde operatör veya kullanıcı başlatma butonuna erken basarsa
sistem oturumu başlatmaz ve kalibrasyon ekranına geri döner. Böylece
araştırma verisine kalibrasyonsuz kayıt karışması önlenir.

\subsection{Final Sahne Ayarları}

Aktif oyun sahnesinde test/build akışı için yapılan temel ayarlar şunlardır:

\begin{itemize}[nosep]
    \item \texttt{GameFlowController.simulatedMode = false},
    \item \texttt{GameFlowController.autoStartAfterCalibration = false},
    \item \texttt{TaskSequencer.autoStartOnAwake = false},
    \item \texttt{LowerLimbBiometrics.useSimulatedInput = false},
    \item \texttt{SessionReportWriter.writeJson = true},
    \item \texttt{ReplayRecorder.autoRecordSession = true}.
\end{itemize}

Ek olarak diagnostics panelinin final kullanıcı akışını karıştırmaması için
otomatik diagnostics kaydı ve panel görünürlüğü kapatılmıştır. Diagnostics
araçları tamamen kaldırılmamıştır; gerektiğinde test amacıyla tekrar
açılabilecek şekilde sahnede korunmuştur.

% ================================================================
\section{Sonuç Doğruluğu İçin Testler ve Veri Yorumlama Güvenliği}
\label{sec:dogrulama}

Kalibrasyon kapısı tamamlandıktan sonra ikinci büyük çalışma, sistemin ürettiği
sonuçların doğru yorumlanabilir olup olmadığının incelenmesidir. Burada amaç,
sistemi tıbbi tanı koyan veya tek başına yaralanma tahmini yapan bir araç gibi
sunmak değildir. Kullanılan ifade çerçevesi özellikle
\textbf{alt ekstremite hareket-risk paterni taraması} olarak sınırlandırılmıştır.

\subsection{JSON ve CSV Raporlarının İncelenmesi}

\texttt{SessionReportWriter} oturum sonunda iki temel çıktı üretir:

\begin{itemize}[nosep]
    \item \texttt{session\_YYYYMMDD\_HHMMSS.csv},
    \item \texttt{session\_YYYYMMDD\_HHMMSS.json}.
\end{itemize}

Bu dosyalarda her görev için süre, valgus, fleksiyon, pelvis salınımı,
simetri indeksi, reach yüzdesi, alt risk bileşenleri, toplam risk, oyun skoru
ve kısa Türkçe yorum saklanmaktadır. Böylece oyun sırasında görülen özet
sonuçlar yalnızca geçici UI metni olarak kalmaz; sonradan tekrar incelenebilen
dosyalar halinde korunur.

\subsection{Üç Tracker Kurulumunda Avatar Kemik Fallback Kullanımı}

Mevcut hardware kurulumunda pelvis ve iki ayak/shin tracker olmak üzere üç
tracker kullanılmaktadır. Bu nedenle diz tracker verisi doğrudan mevcut
değildir. Alt ekstremite açılarının hesaplanabilmesi için
\texttt{LowerLimbBiometrics} içinde avatar IK kemik fallback hattı
kullanılmıştır. Bu sayede kalça-diz-ayak bileği ilişkisi, canlı IK ile sürülen
avatar kemikleri üzerinden çıkarılabilmektedir.

Bu yaklaşım, özellikle diz trackerı bulunmayan build testlerinde süreklilik
sağlar. Bununla birlikte rapor dili, bu verinin bir klinik tanı yerine
tarama ve hareket paterni göstergesi olduğunu açıkça korur.

\subsection{Diagnostics CSV ile Kalite Kontrol}

Kalibrasyon ve açı doğruluğunu denemek için diagnostics CSV kayıtları da
kullanılmıştır. Bu dosyalar örnekleme bazında aşağıdaki alanları içerir:

\begin{itemize}[nosep]
    \item kalibrasyon, IK, tracker ataması ve mapping durumları,
    \item simülasyon girdisi açık/kapalı bilgisi,
    \item sol/sağ bacak veri kullanılabilirliği,
    \item avatar fallback kullanımı,
    \item valgus, fleksiyon, sway, simetri ve reach metrikleri,
    \item pelvis, diz ve ayak bileği pozisyon örnekleri.
\end{itemize}

Bu kayıtlar, kalibrasyonun gerçekten tamamlanıp tamamlanmadığını ve açıların
beklenen büyüklüklerde akıp akmadığını kontrol etmek için kullanılmıştır.
Özellikle önceki manuel marker denemelerinde controller tuş işaretinin güvenilir
olmadığı görülmüş; bu nedenle replay bağlamında otomatik görev olayları
birincil zaman işareti olarak tercih edilmiştir.

% ================================================================
\section{Örnek Oyun Oturumları ve Verilerin Aktarılması}
\label{sec:kayitlar}

Hafta~5 içinde örnek oyun oturumları oynanmış ve headset/build tarafından
üretilen dosyalar bilgisayara aktarılmıştır. Bu dosyalar proje içinde
\texttt{Assets/KayitSonuclari/} klasörü altında toplanmıştır. Böylece veri
analizi yalnızca headset üzerinde değil, Unity Editor ve bilgisayar ortamında
da yapılabilir hale gelmiştir.

Aktarılan kayıt yapısı üç grupta ele alınabilir:

\begin{table}[H]
\centering
\caption{Aktarılan kayıt dosyalarının türleri}
\label{tab:kayit_turleri}
\renewcommand{\arraystretch}{1.22}
\begin{tabularx}{\textwidth}{lX}
\toprule
\textbf{Dosya/Klasör türü} & \textbf{İçerik} \\
\midrule
\texttt{session\_*.json} ve \texttt{session\_*.csv} & Görev bazlı sonuçlar, riskler, skorlar ve yorumlar \\
\texttt{replay\_*/manifest.json} & Replay metadata, oyun zamanı, scene adı, görev listesi, frame/event sayısı \\
\texttt{replay\_*/events.jsonl} & Oturum olayları, görev başlangıç/bitişleri, countdown, rest ve \texttt{result\_ready} olayları \\
\texttt{replay\_*/frames.jsonl} & 30 Hz avatar kemik pozları, kaynak cihaz pozları, IK target pozları ve metrik snapshotları \\
\texttt{lower\_limb\_diagnostics\_*.csv} & Kalibrasyon/tracking/IK kalite kontrol örnekleri \\
\bottomrule
\end{tabularx}
\end{table}

Bu aktarım sürecinin önemli sonucu şudur: oyun verisi artık uygulama kapanınca
kaybolan bir runtime durumundan çıkarılmış, bilgisayarda tekrar okunabilen ve
raporlanabilen bir araştırma çıktısına dönüştürülmüştür.

% ================================================================
\section{Replay Sisteminin Eklenmesi}
\label{sec:replay}

Hafta~5'in en büyük teknik eklerinden biri tam oturum replay sistemidir. Bu
sistem video kaydı üretmez; bunun yerine Unity içinde tekrar oynatılabilen
veri paketi oluşturur. Böylece hem dosya boyutu video kaydına göre daha
yönetilebilir kalır hem de görev adı, yönerge, zaman çizgisi, metrikler ve
sonuç olayları aynı yapı içinde saklanır.

\subsection{Kayıt Tarafı: ReplayRecorder}

\texttt{ReplayRecorder}, aktif oyun sahnesindeki gamification kök objesine
bağlanmıştır. Oturum sequencer üzerinden başladığında replay kaydı otomatik
olarak başlar ve oturum bittiğinde finalize edilir.

Kaydedilen temel veri grupları şunlardır:

\begin{itemize}[nosep]
    \item session ve görev metadata'sı,
    \item Türkçe görev adı ve yönergeler,
    \item countdown, görev başlangıcı, görev bitişi, rest ve session olayları,
    \item görev sonu \texttt{result\_ready} eventleri,
    \item HMD, kontrolcü ve tracker kaynak pozları,
    \item IK target pozları,
    \item final avatar kemik pozları,
    \item canlı metrik snapshotları.
\end{itemize}

Frame kaydı varsayılan olarak 30 Hz yapılmaktadır. Replay paketinin ana
dosyaları manifest, event ve frame kayıtlarıdır:
\begin{center}
\texttt{manifest.json} \quad \texttt{events.jsonl} \quad \texttt{frames.jsonl}
\end{center}

\subsection{Oynatma Tarafı: ReplayPlaybackController ve ReplayAvatarDriver}

Replay oynatma tarafında iki temel bileşen geliştirilmiştir:

\begin{table}[H]
\centering
\caption{Replay oynatma bileşenleri}
\label{tab:replay_components}
\renewcommand{\arraystretch}{1.22}
\begin{tabularx}{\textwidth}{lX}
\toprule
\textbf{Bileşen} & \textbf{Görev} \\
\midrule
\texttt{ReplayPlaybackController} & Replay klasörünü yükler, frame zaman çizgisini oynatır, UI metriklerini günceller \\
\texttt{ReplayAvatarDriver} & Kaydedilmiş avatar kemik pozlarını sahnedeki avatar kemiklerine uygular \\
\texttt{ReplayRecordable} & Gerekirse ek sahne objelerinin transform zincirlerini replay'e dahil eder \\
\bottomrule
\end{tabularx}
\end{table}

Başlangıçta replay sahnesinde \texttt{Animator} bulunmaması durumunda oynatma
aracı avatarı bulamıyordu. Hafta~5 sonunda bu durum düzeltilmiş; sistem
\texttt{SkinnedMeshRenderer.rootBone}, seçili mesh objesi veya doğrudan
\texttt{mixamorig...:Hips} hiyerarşisi üzerinden de kemik eşleştirebilir hale
getirilmiştir.

Ek olarak canlı sahnede \texttt{FullBodyIKSolver} aktif kalırsa, replay'in
uyguladığı kemik pozlarını \texttt{LateUpdate} içinde tekrar bind/IK pozuna
çekebildiği görülmüştür. Bu nedenle replay yüklenirken canlı tracking/IK
sürücüleri kapatılmış ve replay pozlarının uygulanması geç \texttt{LateUpdate}
aşamasına alınmıştır. Overlay üzerinde \texttt{Bones: applied/missing} satırı
ile kemik eşleşme durumu görülebilmektedir.

% ================================================================
\section{Detaylı Analiz Aracının Yazılması}
\label{sec:analiz}

Replay sistemi veri üretimini çözerken, kopyalanan verilerin pratik biçimde
incelenmesi için ayrıca bir Unity Editor analiz aracı yazılmıştır. Bu araçlar
\texttt{ReplayArchiveTools.cs} içinde iki menü olarak sunulmuştur:

\begin{itemize}[nosep]
    \item \texttt{Tools/Gamification/Kayit Analizi},
    \item \texttt{Tools/Gamification/Replay Sahnesi Hazirla}.
\end{itemize}

\subsection{Kayıt Analizi Penceresi}

\texttt{Kayit Analizi} penceresi varsayılan olarak
\texttt{Assets/KayitSonuclari/} klasörünü tarar. \texttt{.meta} dosyalarını
yok sayar ve oturum raporları, replay paketleri ve diagnostics CSV dosyalarını
tek listede gösterir.

Replay ve session raporları birebir aynı zaman damgasına sahip olmayabilir.
Örneğin replay başlangıç zamanı ile session raporunun yazıldığı bitiş zamanı
farklıdır. Bu nedenle araç önce \texttt{manifest.linkedSessionReportJson}
alanını kullanır, ardından timestamp yakınlığı ile eşleştirme yapar.

Pencere içinde aşağıdaki sekmeler bulunmaktadır:

\begin{table}[H]
\centering
\caption{Kayıt analizi penceresi sekmeleri}
\label{tab:analysis_tabs}
\renewcommand{\arraystretch}{1.22}
\begin{tabularx}{\textwidth}{lX}
\toprule
\textbf{Sekme} & \textbf{Gösterilen bilgiler} \\
\midrule
Özet & Oturum adı, süre, görev sayısı, ortalama skor, risk grade dağılımı \\
Görevler & Her görev için skor, risk, valgus, fleksiyon, sway, simetri, reach ve yorum \\
Timeline & \texttt{events.jsonl} içindeki countdown, task start/end, rest ve result eventleri \\
Frame & \texttt{frames.jsonl} üzerinden stream edilen metrik aralıkları ve phase/task frame dağılımı \\
Diagnostics & Kalibrasyon, IK, tracking, mapping sürekliliği ve fallback kullanımı \\
Dosyalar & Replay, manifest, frame, event, session ve diagnostics dosya yolları \\
\bottomrule
\end{tabularx}
\end{table}

Araç ayrıca seçili kayıt için veya tüm kayıtlar için rapor export edebilir.
Çıktılar \path{Assets/KayitSonuclari/AnalizRaporlari/} altında Markdown ve CSV
olarak üretilir.

\subsection{Replay Sahnesini Otomatik Hazırlama}

\texttt{Replay Sahnesi Hazirla} penceresi seçilen replay için sahnedeki gerekli
objeleri otomatik kurar:

\begin{enumerate}[nosep]
    \item Replay klasöründe \texttt{frames.jsonl} varlığını kontrol eder,
    \item Avatar root, mesh veya kemik hiyerarşisini bulur,
    \item \texttt{ReplayAvatarDriver} ekler veya mevcut olanı kullanır,
    \item \texttt{[ReplayReview]} objesini oluşturur,
    \item \texttt{ReplayPlaybackController} ekler ve replay path'ini bağlar,
    \item İstenirse Play Mode'a geçip replay'i otomatik başlatır.
\end{enumerate}

Bu araç, replay inceleme sürecini elle component ekleme ve path bağlama
işlerinden kurtarmaktadır. Böylece headset'ten alınan bir oturum klasörü,
Unity içinde birkaç tıklama ile hem analiz edilebilir hem de görsel olarak
tekrar oynatılabilir hale gelmiştir.

% ================================================================
\section{Hafta 5 Sonunda Oluşan Nihai İş Akışı}
\label{sec:workflow}

Hafta~5 sonunda hedeflenen uçtan uca akış şu hale gelmiştir:

\begin{enumerate}[nosep]
    \item Headset/build açılır ve trackerlar beklenir.
    \item Kullanıcı T-poz kalibrasyonunu tamamlar.
    \item Kalibrasyon, IK ve tracker ataması geçerli değilse oyun başlatılmaz.
    \item Kullanıcı oturumu başlatır.
    \item Görevler sırayla oynanır ve \texttt{TaskEvaluator} her görev sonunda sonuç üretir.
    \item \texttt{SessionReportWriter} JSON/CSV rapor yazar.
    \item \texttt{ReplayRecorder} replay paketini finalize eder.
    \item Dosyalar USB veya dosya aktarımı ile bilgisayara alınır.
    \item \texttt{Kayit Analizi} aracı sonuçları ve frame/timeline verisini inceler.
    \item \texttt{Replay Sahnesi Hazirla} aracı seçili oyunu Unity içinde tekrar oynatır.
\end{enumerate}

Bu akış, projenin araştırma açısından en önemli gereksinimini karşılar:
oyun sırasında üretilen davranış ve skorlar yalnızca anlık UI verisi olarak
kalmaz; kaydedilir, aktarılır, analiz edilir ve gerektiğinde görsel olarak
yeniden oynatılır.

% ================================================================
\section{Hafta 5 Çıktıları ve Değerlendirme}
\label{sec:sonuc}

Hafta~5 sonunda proje, veri güvenilirliği ve sonradan inceleme açısından önemli
bir eşik geçmiştir. Kalibrasyon kapısı sayesinde geçersiz başlangıçlar
engellenmiş, örnek oyunlardan alınan JSON/CSV/replay dosyaları ile kayıt
kalıcılığı doğrulanmış, replay sistemi ile oyun oturumları Unity içinde
yeniden üretilebilir hale getirilmiş ve analiz aracı ile bu veriler daha
okunabilir bir araştırma çıktısına dönüştürülmüştür.

\begin{table}[H]
\centering
\caption{Hafta 5 somut çıktıları}
\label{tab:week5_outputs}
\renewcommand{\arraystretch}{1.22}
\begin{tabularx}{\textwidth}{lX}
\toprule
\textbf{Çıktı} & \textbf{Sağlanan fayda} \\
\midrule
Zorunlu kalibrasyon kapısı & Kalibrasyonsuz veya simüle verili gerçek oturumların başlamasını engelleme \\
JSON/CSV rapor doğrulaması & Görev sonuçlarının oturum sonunda kalıcı ve incelenebilir olması \\
Diagnostics CSV incelemesi & Kalibrasyon, IK, tracking ve metrik sürekliliğini kontrol etme \\
ReplayRecorder & Tam oturumu görev bağlamı, metrikler ve avatar pozlarıyla kaydetme \\
ReplayPlaybackController ve ReplayAvatarDriver & Kaydedilen oyunu Unity içinde headset olmadan oynatma \\
Kayit Analizi aracı & Aktarılan tüm kayıtları listeleme, analiz etme ve rapor export etme \\
Replay Sahnesi Hazırlama aracı & Replay inceleme sahnesini otomatik kurma ve seçili oyunu oynatma \\
\bottomrule
\end{tabularx}
\end{table}

Genel olarak hafta~5, sistemin prototipten araştırma verisi üreten ve ürettiği
veriyi geriye dönük inceletebilen bir yapıya geçiş haftası olmuştur.

\end{document}